% !TeX root = ../../Book3.tex
\section{次调和函数观点}
\label{b3:sec:1803}

Jensen 公式还可以读成一个势论恒等式：圆周平均之差等于圆盘内部 Laplace 测度的对数势。零点只是这一原则的离散实例。

\subsection{\texorpdfstring{\(\log |f|\)}{log|f|}}
\label{b3:subsec:180301}

函数 \(\log|f|\) 在零点外调和，零点处局部为 \(m\log|z-a|\) 加调和函数。这个分解不依赖对数分支选择，因为只取实部。它也使每个局部可积性与奇性符号都可以直接核验。

\subsection{作为 Riesz 测度的零点}
\label{b3:subsec:180302}

\begin{proposition}\label{b3:prop:18-zero-measure}
区域上非恒零全纯函数满足
\[
\frac1{2\pi}\Delta\log|f|=\sum_am_a\,\delta_a,
\]
其中右侧记号 \(m_a\) 表示 \(f\) 在 \(a\) 的零点重数。
\end{proposition}
\begin{proof}
在非零点邻域 Laplace 为零；在零点邻域用局部分解与对数基本解，得到重数乘点质量。零点局部有限，因此各局部恒等式通过有限划分拼成分布等式。
\end{proof}
对亚纯函数，把极点计为负重数，即得相应符号测度。

\subsection{势论解释}
\label{b3:subsec:180303}

若 \(u\in C^2\) 次调和，中心圆周平均 \(m(r)\) 的导数为 \(\mu_u(B(0,r))/r\)，所以
\[
m(R)-u(0)=\int_0^R\frac{\mu_u(B(0,t))}t\,\mathrm{d}t.
\]
证明只需积分圆周平均导数，归一化由 \(\mu_u=\Delta u/(2\pi)\) 确定。光滑逼近可把这一公式推广到适当的次调和代表与非奇中心。对 \(u=\log|f|\)，测度离散，右端正是零点对数计数，因而回到 Jensen。

中心为对数奇点时，\(u(0)=-\infty\)，不能直接相减；必须先分离中心点质量。这是原点重数项出现的原因。

\subsection{与 Nevanlinna 理论的联系}
\label{b3:subsec:180304}

对于亚纯函数，零点与极点可能同时增长。单独看 \(\log|f|\) 的平均会发生抵消，而 \(\log^+|f|\) 的平均与极点积分计数之和能够保留总增长信息。这引向\appref{b3:app:C}的 Nevanlinna 特征。

目前已经得到的是零点的精确势论记录，并没有从 Jensen 一式直接获得第二基本定理。后者还需要对数导数和多个值之间的比较。


\ProseParagraphBreak

一般次调和平均公式可借上一章的局部 Riesz 分解在书内推出。取较大内部圆盘，使 \(u=V+h\)，\(V(z)=\int\log|z-w|\,\mathrm{d}\mu(w)\)，\(h\) 调和。对固定圆周的核平均，有
\[
\frac1{2\pi}\int_0^{2\pi}\log|re^{\mathrm{i}t}-w|\,\mathrm{d}t=\log\max(r,|w|).
\]
若 \(u(0)>-\infty\)，势在中心亦有限，故减去中心值后，核贡献 \(\log^+(r/|w|)\)，这是非负函数。Tonelli 给
\[
\frac1{2\pi}\int u(re^{\mathrm{i}t})\,\mathrm{d}t-u(0)
=\int_{|w|<r}\log(r/|w|)\,\mathrm{d}\mu(w)
=\int_0^r\frac{\mu(B(0,t))}{t}\,\mathrm{d}t.
\]
最后一步将 \(\log(r/|w|)\) 写成 \(\int_{|w|}^r\mathrm{d}t/t\)，再换非负积分次序。圆周上原势的负部可积由此也得到控制；无需未经说明地从分布等式取点值。

\ProseParagraphBreak

这个平均差随 \(r\) 非减，因为 Riesz 测度为正。在对数半径 \(s=\log r\) 下，平均函数的导数为 \(\mu(B(0,e^s))\)，它也非减，所以平均关于 \(s\) 凸。对全纯对数模，导数就是零点整数计数；对一般次调和函数，导数可以连续变化。取 \(u(z)=|z|^2\) 时圆周平均为 \(r^2\)，\(\mu(B(0,r))=2r^2\)，导数为 \(2r\)，恰与公式 \(\mu(B(0,r))/r\) 一致。这同时核验面积归一化和 \(2\pi\) 常数。

\ProseParagraphBreak

若中心是零点，\(u(0)=-\infty\)，则两边不能写作有限数相减。对 \(u=m\log|z|+v\)、\(v(0)\) 有限，圆周平均先保留 \(m\log r\)，再对 \(v\) 使用上述公式；测度中先去掉 \(m\delta_0\)。若任意在原点补一个有限值，Riesz 分布不会因单点修改改变，但次调和平均不等式会失败。这说明分布等价类与点态次调和代表应保持区别，零点处的负无穷值具有数学作用。

\ProseParagraphBreak

对于亚纯函数的 \(\log|f|\)，极点给正无穷的对数奇性，函数本身一般既非次调和也非超调和；它的 Laplace 是符号测度。可以局部写成两个次调和函数之差，对零、极点分别应用平均公式，再取差。符号测度下积分换序需要有限性或正负分解，不能继续把整个表达当成非负积分使用 Tonelli。将正测度论证中的每一步直接复制给符号测度，是从全纯版本过渡到亚纯版本时常见的漏洞。

\ProseParagraphBreak

Nevanlinna 特征保留大小的正部及极点总计数。此处先作一个算术检验：令 \(f(z)=1/z\)，对 \(r>1\)，圆周上 \(\log^+|f|=0\)，但中心极点计数贡献 \(\log r\)，所以特征仍增长；对 \(f(z)=z\)，没有极点，正部平均为 \(\log r\)。两函数互为倒数，粗略增长规模相同，单独的平均 \(\log|f|\) 却一正一负。正部与极点项的结合正是为了防止这种符号抵消丢失增长资料。完整定义及中心归一化将在相应附录中给出，这里只解释为何需要新的量。

\ProseParagraphBreak

零点测度的稳定性也可以从对数模的局部可积收敛看出。下面补上这项分析事实。仍设 \(f_j\to f\not\equiv0\) 局部一致，选一个闭圆盘，其边界无极限零点，内部零点有限。由前面的稳定计数，令 \(p_j\) 为盘中全部零点构成的首一多项式，\(p\) 为极限零点多项式，则 \(p_j\to p\) 一致。商 \(h_j=f_j/p_j\) 在盘内全纯无零点，边界上一致趋于 \(h=f/p\)；最大模原理把收敛传到闭盘。\(h\) 无零点且闭盘上有正下界，故 \(\log|h_j|\to\log|h|\) 一致。
\ProseParagraphBreak
剩下只需处理移动的对数核。若 \(a_j\to a\)，则 \(\log|z-a_j|\to\log|z-a|\) 在固定紧集上依 \(L^1\) 收敛。证明时先删去 \(a\) 附近半径 \(2\delta\) 的小圆，在其外一致收敛；其内，所有充分大的 \(a_j\) 落在半径 \(\delta\) 内，负部积分由
\(2\pi\int_0^{3\delta}r|\log r|\,\mathrm{d}r\) 控制，趋于零。正部在固定有界区域统一有界，所删小圆的面积也趋零。两部分合并即得结论。把有限个移动零点项相加，便得到
\[
 \log|f_j|\longrightarrow\log|f|\quad\text{于 }L^1_{\mathrm{loc}}(D).
\]
对紧支集光滑测试函数 \(\varphi\)，积分 \(\int\log|f_j|\Delta\varphi\) 遂收敛；由零点的 Riesz 测度公式，正好再次得到零点测度的弱收敛。点态对数在移动零点处可能完全没有共同下界，局部可积性才是传递 Laplace 算子的合适层次。
\ProseParagraphBreak
这项结论也说明为什么在正规族中控制函数本身，经常足以控制零点，但前提是先排除恒零极限。若极限恒零，\(\log|f_j|\) 可能整体趋于负无穷，即使零点数并未增加；例如常数 \(1/j\) 的对数在每个点都趋于负无穷，而其 Laplace 始终为零。对数的加性常数与零点测度是两种独立资料，不能从测度收敛反推对数模已经收敛。
